\section{Dirac Probe as Closure-Preserving Solution}
\begin{quote}\small
\noindent\textbf{Status.} \emph{Legacy probe-interface block inside the active main body.}\par
\noindent\textbf{Block function.} This section isolates the Dirac probe as a controlled closure-preserving interface rather than as a new autonomous sector.
\end{quote}
\label{sec:dirac-probe}
%% ============================================================

\begin{remark}[Probe branch as secondary MM/HC transport branch]
The probe sector is not a second ontology beside the main Sphaera line.
It is a secondary branch of the same MM/HC architecture: the same
triolic source, the same closure discipline, and the same compensator
logic are transported into a probe-specific readout regime. The probe
sections should therefore be read as a controlled branch-unfolding of
the common structural core.
\end{remark}

\begin{remark}[The fermionic probe from the tree]
The recursive septonic tree (§\ref{sec:stringlqg}) suggests that
a fermionic probe should not be inserted from outside but should arise
as a special closure-preserving mode.
\end{remark}

\begin{definition}[Meta-Septum closure $C_7$]
The \emph{Meta-Septum closure} is the admissible septonic cluster
selected by a projector $\hat P_{C_7}$ inside $\mathcal{M}_{16}$ such
that
\[
  \mathrm{Tr}\!\bigl(\mathcal{M}_{16}\,\hat P_{C_7}\bigr)
  \equiv 0\pmod{\mathrm{HC}}.
\]
This records the effective closure of the clustered septonic branch
inside the compensator-stabilized shell.
\end{definition}

\begin{theorem}[Dirac probe equation]\label{thm:dirac-probe}
The minimal first-order effective probe equation compatible with Lorentz
covariance, Meta-Septum closure, and the phase-shifted GUE clock
$\hat{\mathcal{R}}_\beta = \hat R_{\mathrm{GUE}}\,e^{i\pi\bsym}$
takes the form
\[
  \bigl(i\gamma^\mu\partial_\mu
        - m(\bsym,\mathrm{GUE})\bigr)\psi = 0,
\]
where $m(\bsym,\mathrm{GUE})$ is the effective inertial scale generated
by the mismatch between the shifted GUE clock and the admissible closure
projector.
Because $\bsym$ is fixed by $\gL^*$ (Theorem~\ref{thm:gue-selfdet}) and
$\gL^*$ is fixed by the Riemann zeros, the mass $m(\bsym,\mathrm{GUE})$
is in principle computable from the zero sequence without any additional
free parameter.
\end{theorem}

\begin{proposition}[Three compatible readings of the fermionic probe]
\label{prop:three-compatible-readings}
Within the present framework, the Dirac probe may be read in three
mutually compatible ways:
\begin{enumerate}[label=(\roman*)]
  \item as a closure-preserving spinorial transport law
        (Theorem~\ref{thm:dirac-probe}),
  \item as a symmetric closure-load response via the anticommutator
        $\langle\psi|\{\hat{\mathcal{R}}_\beta,\hat P_{C_7}\}|\psi\rangle$,
  \item as an antisymmetric shifted-pacer residue via the commutator
        \[
          \mathfrak{m}_{\mathrm{comm}}(\psi)
          = \frac{\hbar}{c}
            \left|\langle\psi|
            [\hat{\mathcal{R}}_\beta,\hat P_{C_7}]|\psi\rangle
            \right|.
        \]
\end{enumerate}
Together: the Dirac probe is a stabilized septonic branch whose
particle-like inertia is generated by the mismatch between admissible
closure and the phase-shifted recursive clock.
The clock scale is fixed by $\gL^*$ (Theorem~\ref{thm:gue-selfdet}),
so these three readings jointly constitute a parameter-free fermionic
probe calculus modulo the explicit computation of
$[\hat{\mathcal{R}}_\beta,\hat P_{C_7}]$.
\end{proposition}


%% ============================================================

